\documentclass[11pt, onecolumn, a4paper]{article}
\usepackage[margin=0.75in]{geometry}
\usepackage{xcolor}
\usepackage{listings}
\usepackage{fancyhdr}
\usepackage{hyperref}
\usepackage{booktabs}
\usepackage{array}

\lstset{
  basicstyle=\ttfamily\small,
  breaklines=true,
  columns=fullflexible,
  backgroundcolor=\color{gray!10},
  frame=single,
  rulecolor=\color{gray}
}

\pagestyle{fancy}
\fancyhf{}
\rhead{\textit{CCEE Exam Prep}}
\lhead{\textbf{Big Data Technologies}}
\cfoot{\thepage}

\title{\textbf{BIG DATA TECHNOLOGIES} \\ \large Comprehensive Cheatsheet for CCEE Exam \\ Sessions 12--19}
\author{Study Guide}
\date{January 2026}

\begin{document}

\maketitle
\tableofcontents
\newpage

\section{SESSION 12: HADOOP ETL DEVELOPMENT}

\subsection{ETL Process Overview}

\textbf{ETL = Extract $\rightarrow$ Transform $\rightarrow$ Load}

\begin{itemize}
  \item \textbf{Extract:} Pull data from OLTP DBs, logs, CSV/JSON files, APIs, message queues $\rightarrow$ HDFS landing zone
  \item \textbf{Transform:} Clean (remove nulls, fix formats), enrich (lookups, joins), aggregate, filter, deduplicate, validate using MapReduce/Hive/Pig/Spark
  \item \textbf{Load:} Store final data into Hive tables, HBase, NoSQL, or downstream RDBMS via connectors (Sqoop)
\end{itemize}

\subsection{ELT (Extract--Load--Transform) Pattern}

Load all raw data into Hadoop data lake first, then build multiple views/data marts later. Common in big data: flexibility to transform data on-demand for different use cases.

\subsection{Typical ETL Functions}

\begin{itemize}
  \item Data cleaning: handle missing values, invalid codes, type conversions
  \item Data integration: joins across sources, key mapping
  \item Data transformation: derived fields (qty $\times$ price), aggregations (SUM, COUNT, AVG, MIN, MAX), sorting, filtering, normalization
  \item Data quality: deduplication, constraint checks, validation rules
\end{itemize}

\subsection{Data Extraction in Hadoop}

\begin{itemize}
  \item From relational DBs: \textbf{Sqoop} or custom connectors/export jobs
  \item From files/logs: \textbf{Flume, Kafka, custom ingestion} to HDFS
  \item From streams: Kafka + Spark/Flink for near real-time ETL
\end{itemize}

\subsection{Why Use ETL Tools (vs hand-coded MapReduce)}

\begin{itemize}
  \item \textbf{Productivity}: Drag--drop or SQL-based design instead of low-level code
  \item \textbf{Reusability \& maintenance}: Centralized job definitions, versioning, documentation
  \item \textbf{Orchestration}: Workflows with dependencies, retries, scheduling, logging, alerts
  \item \textbf{Data quality}: Built-in profiling, validation, error-handling
  \item \textbf{Connectivity}: Ready-made connectors for DBs, SaaS, files, streams, cloud storage
\end{itemize}

\subsection{Advantages of ETL Tools on Hadoop/Cloud}

\begin{itemize}
  \item \textbf{Automation}: Scheduling, monitoring, automatic reruns reduce errors
  \item \textbf{Scalability}: Leverage Hadoop/Spark, scale with cluster size
  \item \textbf{Performance}: Parallel data flows, pushdown optimization, cluster-aware execution
  \item \textbf{Cost-effectiveness}: Commodity hardware + open-source stack; pay-as-you-go cloud models
  \item \textbf{Flexibility}: Support structured, semi-structured, unstructured data (CSV, JSON, logs, images, etc.)
\end{itemize}

\newpage
\section{SESSION 13: HBASE INTRODUCTION}

\subsection{HBase Overview}

\textbf{Distributed, column-oriented NoSQL database} on HDFS. For fast random reads/writes (online access), not batch/stream like HDFS. Ideal for time-series data, clickstreams, IoT metrics, user profiles, random read/write with high throughput.

\subsection{Data Model (Key Terms for MCQs)}

\begin{itemize}
  \item \textbf{Table} $\rightarrow$ rows
  \item \textbf{Row key} $\rightarrow$ sorted lexicographically; unique identifier for each row
  \item \textbf{Column family} $\rightarrow$ logical group of columns; defined at schema level; stored together on disk
  \item \textbf{Column qualifier} $\rightarrow$ specific column inside a family (dynamic, can vary per row)
  \item \textbf{Cell} = (row key, column family, column qualifier, timestamp) $\rightarrow$ value
  \item \textbf{Versions} $\rightarrow$ identified by timestamp; each cell can have multiple timestamped versions
\end{itemize}

\subsection{HBase Architecture}

\begin{itemize}
  \item \textbf{HMaster}: Lightweight master managing cluster, assigning regions to RegionServers, DDL (create/alter/delete), load balancing
  \item \textbf{RegionServer (HRegionServer)}: Worker serving read/write requests for assigned regions
  \begin{itemize}
    \item \textbf{MemStore}: Write cache in RAM
    \item \textbf{HFile}: On-disk store in HDFS
    \item \textbf{BlockCache}: Read cache
    \item \textbf{WAL (Write-Ahead Log)}: Durability log
  \end{itemize}
  \item \textbf{ZooKeeper}: Coordination service for master election, server registration, metadata discovery by clients
  \item \textbf{Regions}: Horizontal partitions (row-key ranges) of a table
  \begin{itemize}
    \item Each region served by one RegionServer at a time
    \item Automatically split when too large for load distribution
  \end{itemize}
\end{itemize}

\subsection{Installation Conceptual Steps}

\begin{enumerate}
  \item \textbf{Prerequisites}: Working Hadoop cluster (HDFS), ZooKeeper ensemble
  \item \textbf{Configuration}: Download HBase, configure \texttt{hbase-site.xml} pointing to HDFS and ZooKeeper
  \item \textbf{Start services}: Start ZooKeeper, HMaster, RegionServers
  \item \textbf{Verify}: Use HBase shell (\texttt{create}, \texttt{list}, \texttt{put}, \texttt{get}, \texttt{scan}) to verify
\end{enumerate}

\subsection{Installation Modes}

\begin{itemize}
  \item \textbf{Standalone}: Single JVM, local file system (dev only)
  \item \textbf{Pseudo-distributed}: All daemons on one machine with HDFS (testing)
  \item \textbf{Fully distributed}: Multi-node cluster with HDFS and ZooKeeper (production)
\end{itemize}

\newpage
\section{SESSION 14: HBASE ADMIN, OPERATIONS \& SECURITY}

\subsection{HBase Shell: Administrative (DDL) Commands}

\begin{itemize}
  \item \texttt{create 't1', 'cf1','cf2'} -- Create table with column families
  \item \texttt{disable 't1'}, \texttt{enable 't1'} -- Change table state
  \item \texttt{alter 't1', \{NAME => 'cf1', VERSIONS => 3\}} -- Modify schema
  \item \texttt{describe 't1'} -- Show table schema
  \item \texttt{drop 't1'} -- Delete table (removes metadata and data)
\end{itemize}

\subsection{HBase Shell: Data (CRUD/DML) Commands}

\begin{itemize}
  \item \texttt{put 't1','row1','cf1:col1','value1'} -- Create/Update
  \item \texttt{get 't1','row1'} -- Read one row
  \item \texttt{scan 't1'} -- Read many rows (can specify filters, columns, LIMIT)
  \item \texttt{delete 't1','row1','cf1:col1'} or \texttt{deleteall 't1','row1'} -- Delete
  \item \texttt{count 't1'} -- Count rows
  \item \texttt{truncate 't1'} -- Disable, drop, recreate table (clears data, keeps structure)
\end{itemize}

\subsection{Java Client API (org.apache.hadoop.hbase.client) -- DML/CRUD}

\begin{itemize}
  \item \texttt{Connection} / \texttt{ConnectionFactory} -- Create cluster connection
  \item \texttt{Table} -- Logical handle to HBase table
  \item \texttt{Put} -- Create/update row
  \item \texttt{Get} -- Retrieve row or specific columns
  \item \texttt{Delete} -- Delete columns/rows
  \item \texttt{Scan} -- Iterate over ranges of rows
  \item \texttt{Result} -- Represents data returned from Get/Scan
\end{itemize}

\subsection{HBase Admin API (HBaseAdmin / Admin) -- DDL}

\begin{itemize}
  \item \texttt{admin.createTable(...)} -- Create table
  \item Disable/enable table
  \item Add/modify/delete column families
  \item Get table descriptors, list tables
\end{itemize}

\subsection{HBase Security}

\begin{itemize}
  \item \textbf{Authentication}: \textbf{Kerberos-based} with SASL at RPC level; only authenticated users/clients can communicate with HBase
  \item \textbf{Authorization}: Implemented via \textbf{Access Controller / ACL coprocessor}
  \begin{itemize}
    \item Grants permissions: R (Read), W (Write), X (Execute), C (Create), A (Admin) at table/family/qualifier level
    \item \texttt{grant 'user','RWXCA','table','cf','col'} -- Grant permissions
    \item \texttt{revoke 'user','table'} -- Remove permissions
  \end{itemize}
\end{itemize}

\newpage
\section{SESSION 16: HIVE DATA WAREHOUSE -- INTRODUCTION}

\subsection{Hive Overview}

\textbf{Data warehouse system on Hadoop} using SQL-like queries (HiveQL) over HDFS. For \textbf{batch analytics} (OLAP), not OLTP. Created at Facebook to let analysts write SQL over massive log data.

\subsection{Hive Architecture Components}

\begin{itemize}
  \item \textbf{HiveQL}: SQL-like query language
  \item \textbf{Driver}: Receives query, manages session, coordinates execution (parser + planner + optimizer)
  \item \textbf{Compiler}: Parses query, performs semantic analysis, generates execution plan
  \item \textbf{Metastore}: RDBMS-backed metadata store (table schemas, partitions, locations, SerDe info)
  \item \textbf{Execution engine}: Converts plan to jobs on Hadoop: MapReduce / Tez / Spark
  \item \textbf{HiveServer2 / Clients}: JDBC/ODBC/CLI interfaces for apps and users
\end{itemize}

\subsection{Hive Installation (Conceptual)}

\begin{enumerate}
  \item Install Hadoop (HDFS/YARN)
  \item Install Hive binaries
  \item Configure metastore DB (MySQL/PostgreSQL/Derby) and \texttt{hive-site.xml}
  \item Start metastore service and HiveServer2
  \item Connect via CLI or Beeline; run HiveQL queries
\end{enumerate}

\subsection{Hive vs Traditional RDBMS}

\begin{center}
\begin{tabular}{|l|p{3cm}|p{3cm}|}
\hline
\textbf{Aspect} & \textbf{HIVE} & \textbf{RDBMS} \\
\hline
Workload & OLAP/batch; scans huge datasets & OLTP; many small reads/writes \\
\hline
Latency & High (batch); startup cost & Low; interactive \\
\hline
Storage & HDFS; horizontal scale to TB--PB & Local/RAID; vertical scale; indexes \\
\hline
Schema & Schema-on-read; data dumped first & Schema-on-write; schema enforced at insert \\
\hline
Execution & Distributed jobs (MR/Tez/Spark) & Single server/cluster, low-latency engine \\
\hline
Transactions & No strict row-level ACID & Strict ACID by default \\
\hline
\end{tabular}
\end{center}

\newpage
\section{SESSION 17: WORKING WITH HIVE QL}

\subsection{Datatypes}

\textbf{Primitive Types:}
\begin{itemize}
  \item Numeric: TINYINT, SMALLINT, INT, BIGINT, FLOAT, DOUBLE, DECIMAL
  \item String: STRING, VARCHAR, CHAR
  \item Date/time: DATE, TIMESTAMP
  \item Boolean: BOOLEAN
\end{itemize}

\textbf{Complex Types:}
\begin{itemize}
  \item ARRAY$<$type$>$
  \item MAP$<$key\_type, value\_type$>$
  \item STRUCT$<$field1:type1, field2:type2, ...$>$
  \item UNIONTYPE
\end{itemize}

\subsection{Operators \& Built-in Functions}

\begin{itemize}
  \item \textbf{Relational}: =, !=, $<$, $<$=, $>$, $>$=
  \item \textbf{Arithmetic}: +, -, *, /, \%
  \item \textbf{Logical}: AND, OR, NOT
  \item \textbf{String}: LIKE, RLIKE, CONCAT, SUBSTR, LENGTH, UPPER, LOWER, TRIM
  \item \textbf{Aggregate}: SUM, COUNT, AVG, MIN, MAX, STDDEV
  \item \textbf{Math}: FLOOR, CEIL, ROUND, ABS, SQRT
\end{itemize}

\subsection{Managed (Internal) Tables}

\begin{itemize}
  \item Created by default (no \texttt{EXTERNAL} keyword)
  \item Hive \textbf{owns the data}; stored under warehouse path (e.g., \texttt{/user/hive/warehouse/dbname.db/tablename/})
  \item \texttt{DROP TABLE} deletes \textbf{both metadata and data} from HDFS
\end{itemize}

\subsection{External Tables}

\begin{itemize}
  \item Created with \texttt{CREATE EXTERNAL TABLE ... LOCATION 'path'}
  \item Hive manages \textbf{only metadata}; data resides in specified location
  \item \texttt{DROP TABLE} by default removes \textbf{only metadata}, not actual files
  \item Good for sharing data or data managed by external systems
\end{itemize}

\subsection{Partitions}

\begin{itemize}
  \item Directory-level split on partition columns
  \item Created with \texttt{PARTITIONED BY (dt STRING, region STRING)}
  \item Data physically stored in subdirectories: \texttt{.../table/dt=2025-01-01/region=IN/}
  \item \textbf{Partition pruning}: queries filtering on partition columns can skip other partitions $\rightarrow$ major performance boost
\end{itemize}

\subsection{Buckets}

\begin{itemize}
  \item File-level split inside partitions
  \item Created with \texttt{CLUSTERED BY (user\_id) INTO 8 BUCKETS}
  \item Data hashed on bucket column, written to multiple bucket files
  \item Helps with sampling, map-side joins, query optimization
\end{itemize}

\subsection{Storage Formats}

\begin{itemize}
  \item \textbf{Row-oriented}: TEXTFILE, AVRO (less efficient for analytical scans)
  \item \textbf{Columnar}: ORC, PARQUET (better compression, faster analytical queries)
  \item Columnar formats preferred for modern Hive analytics
\end{itemize}

\subsection{Importing Data (DML Operations)}

\begin{itemize}
  \item \texttt{LOAD DATA INPATH '/path/file.csv' INTO TABLE t;} -- Move/copy files from HDFS into table location
  \item \texttt{LOAD DATA LOCAL INPATH 'localfile' INTO TABLE t;} -- From local file system
  \item External sources via Sqoop, Spark, or ETL tools
\end{itemize}

\subsection{Altering Tables (DDL)}

\begin{itemize}
  \item Add partition: \texttt{ALTER TABLE t ADD PARTITION (dt='2025-01-01');}
  \item Add/replace columns: \texttt{ALTER TABLE t ADD COLUMNS (new\_col STRING);}
  \item Change properties: \texttt{ALTER TABLE t SET TBLPROPERTIES (...);}
  \item Change file format: \texttt{ALTER TABLE t SET FILEFORMAT ...;}
\end{itemize}

\subsection{Dropping Tables}

\texttt{DROP TABLE t;}
\begin{itemize}
  \item \textbf{Managed table}: drops metadata + data
  \item \textbf{External table}: drops metadata only; data files remain
\end{itemize}

\newpage
\section{SESSION 18: QUERYING WITH HIVE QL}

\subsection{Sorting Clauses}

\begin{itemize}
  \item \texttt{ORDER BY} -- \textbf{Global sort} using a single reducer; entire result is totally ordered
  \item \texttt{SORT BY} -- Sorts data \textbf{per reducer only}; overall output may not be globally sorted
  \item \texttt{DISTRIBUTE BY} -- Controls which rows go to which reducer (by key)
  \item \texttt{CLUSTER BY} -- Shorthand for \texttt{DISTRIBUTE BY} + \texttt{SORT BY} on same columns
\end{itemize}

\subsection{Aggregation with GROUP BY}

\begin{itemize}
  \item Standard aggregate functions: SUM, COUNT, AVG, MIN, MAX, STDDEV
  \item \texttt{GROUP BY} groups rows by specified columns
  \item \texttt{HAVING} filters \textbf{groups after aggregation} (vs \texttt{WHERE} filters rows before grouping)
\end{itemize}

\subsection{Joins in HiveQL}

\begin{itemize}
  \item Supported types: INNER, LEFT OUTER, RIGHT OUTER, FULL OUTER, CROSS, SEMI
  \item Basic syntax: \texttt{SELECT * FROM a JOIN b ON a.id = b.id;}
  \item \textbf{Reduce-side (common) join / shuffle join}:
  \begin{itemize}
    \item All data for a join key shuffled to the same reducer, joined there
    \item Works for arbitrary-sized tables; involves heavy network shuffle
  \end{itemize}
  \item \textbf{Map-side join (MapJoin)}:
  \begin{itemize}
    \item One table must be small enough to fit in mapper memory
    \item Small table loaded on each mapper; mappers stream big table and join locally
    \item No reduce phase; less shuffle $\rightarrow$ better performance
    \item Enabled with hint: \texttt{SELECT /*+ MAPJOIN(small\_table) */ ...}
  \end{itemize}
\end{itemize}

\subsection{Subqueries}

\begin{itemize}
  \item \textbf{Subqueries in FROM clause} (derived tables):
  \begin{itemize}
    \item Pre-aggregate data before joining
    \item Calculate top-N per group, then filter on that result
  \end{itemize}
  \item \textbf{Subqueries in WHERE clause}:
  \begin{itemize}
    \item Scalar subqueries for filtering based on another query result
  \end{itemize}
  \item Combine filtering + aggregation in single logical step
\end{itemize}

\subsection{Views}

\begin{itemize}
  \item \textbf{Virtual table} defined by a saved query; no separate data storage
  \item \texttt{CREATE VIEW top\_regions AS SELECT region, SUM(sales) s FROM sales GROUP BY region;}
  \item Query like a table; underlying base tables are scanned
  \item Benefits: centralize logic; changes apply to all consumers; simpler interface for complex queries
\end{itemize}

\subsection{MapReduce Scripts (TRANSFORM Clause)}

\begin{itemize}
  \item Integrate external scripts in MapReduce-like fashion
  \item Syntax: \texttt{SELECT TRANSFORM(col1, col2) USING 'myscript.py' AS (out1, out2) FROM t;}
  \item Script reads tab-delimited input from STDIN, writes output to STDOUT
  \item Hive handles parallelization and MapReduce plumbing
  \item Older syntax: \texttt{MAP ... REDUCE ...} clauses
\end{itemize}

\newpage
\section{SESSION 19: MORE ON HIVE QL (DATA MANIPULATION, UDFS, SCRIPTS)}

\subsection{Data Manipulation (DML)}

\begin{itemize}
  \item \texttt{INSERT INTO TABLE t SELECT ...} -- \textbf{Append} data to existing table/partition
  \item \texttt{INSERT OVERWRITE TABLE t SELECT ...} -- \textbf{Replace} existing data in table/partition
  \item \texttt{LOAD DATA INPATH 'hdfs\_path' INTO TABLE t} -- Move/copy files into table (append at file level)
  \item HDFS is write-once; updates/deletes done by rewriting partitions with \texttt{INSERT OVERWRITE}
\end{itemize}

\subsection{User-Defined Functions (UDFs)}

\textbf{Types:}
\begin{itemize}
  \item \textbf{UDF}: Scalar function (row-by-row transformation)
  \item \textbf{UDAF}: Aggregate function (e.g., custom SUM-like behavior)
  \item \textbf{UDTF}: Table-generating function (expands one row into multiple rows)
\end{itemize}

\textbf{Registration:}
\begin{enumerate}
  \item \texttt{ADD JAR /path/myudf.jar;}
  \item \texttt{CREATE TEMPORARY FUNCTION myfunc AS 'com.example.MyUdf';}
  \begin{itemize}
    \item Temporary functions live for the session only
    \item Permanent functions stored in metastore (version-dependent)
  \end{itemize}
\end{enumerate}

\subsection{Custom MapReduce in Hive}

\begin{itemize}
  \item Via \texttt{TRANSFORM}, \texttt{MAP}, \texttt{REDUCE} clauses
  \item Script acts as mapper/reducer: reads tab-delimited STDIN, writes tab-delimited STDOUT
  \item Hive distributes and manages MapReduce execution
  \item Allows any language (Python, Bash, etc.) for complex transformations
\end{itemize}

\subsection{Appending Data into Existing Tables}

\begin{itemize}
  \item \texttt{INSERT INTO TABLE sales SELECT ...;} -- Append into table
  \item \texttt{INSERT INTO TABLE sales\_p PARTITION (dt='2026-01-05') SELECT ...;} -- Append into partition
  \item \texttt{LOAD DATA INPATH '/data/file' INTO TABLE sales;} -- Append files
\end{itemize}

\subsection{Writing HQL Scripts (.hql Files)}

\begin{itemize}
  \item Put sequence of HiveQL commands in \texttt{.hql} file (each ending with \texttt{;})
  \item Execute from shell:
  \begin{itemize}
    \item \texttt{hive -f myscript.hql}
    \item \texttt{beeline -f myscript.hql}
  \end{itemize}
  \item Used to automate ETL jobs, daily reports, multi-step operations
  \item Provides reusable, version-controlled pipeline definitions
  \item Commonly scheduled in production via Oozie, Airflow, or other schedulers
\end{itemize}

\subsection{Common HiveQL SELECT Clauses (Summary)}

\begin{lstlisting}[language=SQL]
SELECT columns
FROM table
WHERE row_filter
GROUP BY group_columns
HAVING group_filter
ORDER BY / SORT BY / DISTRIBUTE BY / CLUSTER BY columns
LIMIT count;
\end{lstlisting}

\newpage
\section{KEY MCQ TRAPS \& MNEMONICS}

\subsection{ETL vs ELT}

\begin{itemize}
  \item \textbf{ETL}: Extract $\rightarrow$ Transform $\rightarrow$ Load (traditional DWH)
  \item \textbf{ELT}: Extract $\rightarrow$ Load $\rightarrow$ Transform (modern data lakes; load all raw data first, transform on-demand)
\end{itemize}

\subsection{Hive vs RDBMS in One Sentence}

\begin{itemize}
  \item \textbf{Hive}: ``SQL queries on huge HDFS files executed as distributed batch jobs''
  \item \textbf{RDBMS}: ``Row-level transactions on indexed local storage at interactive speeds''
\end{itemize}

\subsection{Managed vs External Table on DROP}

\begin{itemize}
  \item \textbf{Managed}: \texttt{DROP TABLE} removes metadata \textbf{and} data $\rightarrow$ data gone forever
  \item \textbf{External}: \texttt{DROP TABLE} removes metadata only $\rightarrow$ data files remain accessible
\end{itemize}

\subsection{Partitions vs Buckets}

\begin{itemize}
  \item \textbf{Partitions}: Directory-level split $\rightarrow$ prune entire directories in queries $\rightarrow$ huge performance gains
  \item \textbf{Buckets}: File-level split within partitions $\rightarrow$ helps joins/sampling, not query pruning
\end{itemize}

\subsection{Reduce-side vs Map-side Join}

\begin{itemize}
  \item \textbf{Reduce-side}: Any table size; heavy shuffle (default)
  \item \textbf{Map-side}: One table must fit mapper memory; minimal shuffle (faster)
\end{itemize}

\subsection{Sorting Clauses Ranking (Performance Impact)}

\begin{itemize}
  \item Fastest: \texttt{SORT BY} (per-reducer sort only)
  \item Medium: \texttt{DISTRIBUTE BY} (partition data, no sort)
  \item Slowest: \texttt{ORDER BY} (global sort, single reducer bottleneck)
\end{itemize}

\subsection{Hive Operations Cheat}

\begin{itemize}
  \item \textbf{DDL (schema)}: CREATE TABLE, ALTER TABLE, DROP TABLE
  \item \textbf{DML (data)}: INSERT INTO, INSERT OVERWRITE, LOAD DATA
  \item \textbf{Query}: SELECT with WHERE/GROUP BY/HAVING/ORDER BY
  \item \textbf{UDF}: Custom Java logic registered via ADD JAR + CREATE FUNCTION
\end{itemize}

\newpage
\section{QUICK REFERENCE: ESSENTIAL COMMANDS}

\subsection{HBase Shell (Most Common)}

\begin{lstlisting}
hbase(main):001:0> create 't1', 'cf1'
hbase(main):002:0> put 't1', 'row1', 'cf1:col1', 'value1'
hbase(main):003:0> get 't1', 'row1'
hbase(main):004:0> scan 't1'
hbase(main):005:0> count 't1'
hbase(main):006:0> disable 't1' -> drop 't1'
\end{lstlisting}

\subsection{Hive QL (Most Common)}

\begin{lstlisting}[language=SQL]
CREATE TABLE sales (id INT, region STRING, amount DOUBLE);
LOAD DATA INPATH '/data/sales.csv' INTO TABLE sales;
SELECT region, SUM(amount) FROM sales WHERE amount > 0 GROUP BY region;
INSERT OVERWRITE TABLE sales_summary 
  SELECT region, SUM(amount) FROM sales GROUP BY region;
CREATE VIEW high_sales AS 
  SELECT region, amount FROM sales WHERE amount > 5000;
\end{lstlisting}

\subsection{Hive MapReduce Script Integration}

\begin{lstlisting}[language=SQL]
SELECT TRANSFORM(col1, col2)
USING 'myscript.py'
AS (output1, output2)
FROM input_table;
\end{lstlisting}

\subsection{Execute HQL Script}

\begin{lstlisting}[language=bash]
$ hive -f queries.hql
$ beeline -f queries.hql
\end{lstlisting}

\newpage
\section{EXAM STRATEGY REMINDERS}

\begin{enumerate}
  \item \textbf{Always distinguish}:
  \begin{itemize}
    \item ETL (traditional) vs ELT (data lake pattern)
    \item Managed vs External tables (especially DROP behavior)
    \item Reduce-side vs Map-side joins (default vs optimized)
  \end{itemize}

  \item \textbf{Data model clarity}:
  \begin{itemize}
    \item Row keys are sorted lexicographically (HBase)
    \item Partitions are directories (Hive)
    \item Buckets are files within partitions (Hive)
  \end{itemize}

  \item \textbf{Scaling rules}:
  \begin{itemize}
    \item HDFS scales via replication + rack awareness
    \item HBase scales via region splitting
    \item Hive scales via partitions (prune), buckets, and distributed execution engine
  \end{itemize}

  \item \textbf{Performance patterns}:
  \begin{itemize}
    \item Map-side join $>$ Reduce-side join (when small table fits memory)
    \item Column-oriented formats (ORC, PARQUET) $>$ row-oriented (TEXT) for analytics
    \item Partition pruning filters before shuffle (huge savings)
  \end{itemize}

  \item \textbf{MCQ red flags}:
  \begin{itemize}
    \item ``Always'' or ``never'' statements are often wrong
    \item Details matter: \texttt{DROP TABLE} on managed vs external, \texttt{INSERT INTO} vs \texttt{INSERT OVERWRITE}
    \item Multiple-correct-answer questions are common
  \end{itemize}
\end{enumerate}

\newpage
\section{FINAL SUMMARY TABLE}

\begin{center}
\begin{tabular}{|l|l|p{5cm}|}
\hline
\textbf{Component} & \textbf{Core Purpose} & \textbf{Key Concepts} \\
\hline
ETL & Data pipeline design & Extract, Transform, Load; ELT pattern; Sqoop, Flume, Kafka \\
\hline
HBase & Fast random read/write NoSQL & Row keys, column families, regions, RegionServers, WAL, ZooKeeper \\
\hline
Hive & Batch SQL analytics on HDFS & Metastore, partitions, buckets, managed/external, MapReduce execution \\
\hline
Joins & Combine multiple data sources & Reduce-side (default), map-side (small table in memory) \\
\hline
UDFs & Custom logic & Java JARs, ADD JAR, CREATE TEMPORARY FUNCTION \\
\hline
Scripts & Automation & HQL files, \texttt{hive -f}, orchestration scheduling \\
\hline
\end{tabular}
\end{center}

\vspace{1cm}

\begin{center}
\textbf{Good luck on your CCEE exam!} \\
Study the MCQ patterns, understand the core concepts deeply, \\
and remember the key differences between similar-sounding terms.
\end{center}

\end{document}